\section{Svar på spørsmål}

\subsection{Hvordan (eller kan) dere sørge for at en instruktør ikke kan være til stede på to
forskjellige plasser samtidig? Kan dere få til dette ved hjelp av modellen deres,
eller må noe slikt implementeres i programvaren}

I vår modell lagres instruktøren for en gruppeaktivitet som fremmednøkkelen
\textit{instrukt\_ID} i \textit{Gruppeaktivitet}-tabellen.
Det finnes ingen restriksjon i databaseskjemaet som hindrer at samme
\textit{instrukt\_ID} opptrer i to gruppeaktiviteter med overlappende
tidsrom på samme dato.

En \texttt{UNIQUE}-restriksjon kan kun forhindre identiske verdier, ikke
overlappende \emph{intervaller}.
To aktiviteter for samme instruktør på samme dato med delvis
overlappende start-/slutt-tidspunkt ville derfor ikke bli fanget opp.

\paragraph{Konklusjon.}
Denne regelen \emph{kan ikke} uttrykkes i den relasjonelle modellen
og må derfor implementeres i applikasjonskoden (eller via en trigger).
Før en gruppeaktivitet opprettes eller oppdateres, må applikasjonen
sjekke at det ikke allerede finnes en annen aktivitet med samme
instruktør der tidsintervallet overlapper.

\subsection{Det samme spørsmålet for en bruker. Kan dette løses i modellen eller må det
gjøres i programvaren?}

Resonnementet er det samme som for instruktører, men situasjonen er
noe bredere fordi en bruker kan delta i flere typer aktiviteter.
En profil kan være registrert i \textit{påmeldt\_til} for flere
gruppeaktiviteter, og i \textit{møter\_til\_idrett} for flere
idrettslagstimer.
Det finnes ingen restriksjon i skjemaet som forhindrer overlappende
tidsrom, verken innenfor samme aktivitetstype eller på tvers av de to.
En bruker kan altså potensielt kollidere på tre måter:
\begin{enumerate}
    \item to gruppeaktiviteter som overlapper,
    \item to idrettslagstimer som overlapper,
    \item en gruppeaktivitet og en idrettslagstime som overlapper.
\end{enumerate}

I alle tilfellene krever overlappsjekken sammenligning av
\textit{dato}, \textit{start} og \textit{slutt} på tvers av rader
(og potensielt på tvers av tabellene \textit{Gruppeaktivitet} og
\textit{Idrettslagstime}) for samme \textit{profil\_ID}.
Siden tidsattributtene ikke ligger i relasjonstabellene
\textit{påmeldt\_til} og \textit{møter\_til\_idrett}, men i
\textit{Gruppeaktivitet} og \textit{Idrettslagstime}, kreves det
en også en join for å utføre sjekken.
Dette kan ikke uttrykkes med \texttt{UNIQUE} eller \texttt{CHECK}.

\paragraph{Konklusjon.}
Også denne regelen må håndteres i applikasjonskoden.
Ved påmelding til en gruppeaktivitet eller oppmøte til en
idrettslagstime bør applikasjonen kontrollere at brukeren ikke
allerede deltar i en annen aktivitet som overlapper i tid.

\subsection{Fra hvilket usecase testes det for / opprettes utestengelse (svartelisting)? Det
trenger ikke være en av de oppgitte usecasene.}

Utestengelse involverer to steg: (1) \emph{oppretting} av prikker som
til slutt utløser utestengelsen, og (2) \emph{sjekk} av om en profil
er utestengt.

\paragraph{Oppretting.}
Prikker gis til profiler som er påmeldt en gruppeaktivitet uten å
registrere oppmøte senest 5~minutter før start.
Grunnlaget for dette er \textbf{brukstilfelle~3} (registrering av
oppmøte): det er oppmøteregistreringene i \textit{møter\_til\_gruppe}
som gjør det mulig å avgjøre hvem som \emph{ikke} har møtt.
Selve tildelingen av prikker skjer i et eget brukstilfelle;
\emph{tildeling av prikker etter oppmøtefrist},
som etter at oppmøtefristen har passert (5~minutter før start),
sammenligner \textit{påmeldt\_til} med \textit{møter\_til\_gruppe}
og oppretter en rad i \textit{Prikk} for de som ikke har møtt.
Når en profil har fått $\geq 3$ prikker innenfor de siste 30 dagene,
er den utestengt.

\paragraph{Sjekk.}
Utestengelsessjekken er implementert direkte i databasen ved hjelp av
to triggere på tabellen \textit{påmeldt\_til}:
\begin{itemize}
    \item \texttt{check\_utestengelse\_insert} – kjøres \texttt{BEFORE INSERT}
    \item \texttt{check\_utestengelse\_update} – kjøres \texttt{BEFORE UPDATE}
\end{itemize}
Dersom profilen har $\geq 3$ prikker de siste 30 dagene, avbrytes
operasjonen med \texttt{RAISE(ABORT, \ldots)}.
Sjekken trenger derfor ikke implementeres i applikasjonslaget
– databasen håndhever regelen automatisk ved hvert forsøk på påmelding.

Brukstilfelle~6 demonstrerer svartelistingsmekanismen eksplisitt, men
den naturlige triggeren i et produksjonssystem er brukstilfellet for
prikk-tildeling beskrevet over, kombinert med triggerne i databasen
som blokkerer påmelding ved utestengelse.

\subsection{Statistikk er det noe som må lagres eksplisitt eller går det an å bare gjøre queries
for å få svar? Diskuter dette.}

Vi mener det meste av statistikken som er relevant for SiT kan utledes med vanlige
SQL-spørringer mot databasen. Dette vil også gjøre det enklere å oppdatere statistikken hvis det skjer endringer i databasen.

\paragraph{Eksempler på statistikk som kan beregnes med queries.}
\begin{itemize}
    \item \textbf{Antall påmeldte per økt:}
          \texttt{COUNT(*)} på \textit{påmeldt\_til} gruppert på gruppeaktivitet.
    \item \textbf{Hvem som deltok på en gitt trening:}
          spørring mot \textit{møter\_til\_gruppe}.
    \item \textbf{Personen(e) med flest gruppetimer i en måned} (brukstilfelle~7):
          \texttt{COUNT(*)} på \textit{møter\_til\_gruppe} gruppert på
          \textit{profil\_ID}, filtrert på måned.
    \item \textbf{Par som trener sammen} (brukstilfelle~8):
          selvjoin på \textit{møter\_til\_gruppe}.
    \item \textbf{Ukeplan} (brukstilfelle~4):
          spørring mot \textit{Gruppeaktivitet} filtrert på datoperiode.
    \item \textbf{Prikkstatus og utestengelse:}
          beregnes fra \textit{Prikk}-tabellen ved å telle prikker
          innenfor de siste 30 dagene.
\end{itemize}

\paragraph{Konklusjon.}
Statistikk trenger \emph{ikke} lagres eksplisitt i vår modell;
alt kan utledes av de normaliserte tabellene med vanlige SQL-spørringer.
